\documentclass[11pt,a4paper]{article}
\usepackage[left=2cm,right=2cm,top=2.5cm,bottom=2.5cm]{geometry}
\usepackage{fontspec}
\usepackage{ctex}
\usepackage{amsmath,amssymb,amsfonts,mathtools}
\usepackage[mathbf=sym]{unicode-math}
\setmainfont{STIX Two Text}
\setmathfont{STIX Two Math}
\setsansfont{Arial}
\setmonofont{Courier New}
\setCJKmainfont[BoldFont=Noto Serif CJK SC Bold]{Noto Serif CJK SC}
\setCJKsansfont[BoldFont=Noto Sans CJK SC Bold]{Noto Sans CJK SC}
\setCJKmonofont{Noto Sans CJK SC}
\newCJKfontfamily\examkaishu{FandolKai-Regular}
\usepackage{graphicx}
\usepackage{booktabs,array,multirow,tabularx}
\usepackage{fancyhdr}
\usepackage{needspace}
\usepackage{tikz}
\usepackage{hyperref}
\hypersetup{colorlinks=false}
\usepackage{parskip}
\setlength{\parskip}{0.4em}
\setlength{\parindent}{0em}
\usepackage{xcolor}

% 数学符号
\renewcommand{\ge}{\geqslant}
\renewcommand{\geq}{\geqslant}
\renewcommand{\le}{\leqslant}
\renewcommand{\leq}{\leqslant}

% 知识点标签
\newcommand{\kp}[1]{%
    \par\vspace{0.3em}\noindent%
    {\small\textcolor{gray}{#1}}%
    \par\vspace{0.2em}%
}

% 答案解析区域
\newenvironment{solutionblock}{%
    \par\vspace{0.5em}%
    \noindent\rule{\textwidth}{0.4pt}%
    \vspace{0.5em}%
    \noindent\textbf{【参考答案与解析】}%
    \par\vspace{0.3em}%
    {\small%
}{%
    \par}%
}

% 答案标签
\newcommand{\anslabel}[1]{%
    \par\vspace{0.2em}\noindent%
    \textbf{#1}%
    \ignorespaces%
}

% 页眉页脚
\pagestyle{fancy}
\fancyhf{}
\fancyhead[L]{\small\examkaishu 虹口区高三二模数学（教师版）}
\fancyhead[R]{\small\examkaishu 第 \thepage 页}
\fancyfoot[C]{}
\renewcommand{\headrulewidth}{0.4pt}

\begin{document}

{\centering\Large\bfseries 2025-2026学年虹口区高三二模数学试卷（教师版）\par}
\vspace{0.3em}
\hrule
\vspace{0.5em}

# 2025-2026学年虹口区高三二模数学试卷（教师版）

\textbf{含参考答案及详细解析}
---


## 第 1 题


1. 设全集为 $U = \{  - 2, - 1,0,1,2,3\}$ ，集合 $A = \left\{  {x\left| {\;\frac{x}{x - 2} \geq  0}\right. , x \in  U}\right\}$ ，则 $\bar{A} =$ ___.



\begin{solutionblock}
{\small \anslabel{【考点】：} 集合运算、分式不等式}
{\small \anslabel{【答案】} $\{ 1,2\}$}
{\small \anslabel{【解析】}}
{\small \anslabel{【分析】} 根据分式不等式的解法,可求得集合 $A$ ,即可得答案.}
{\small \anslabel{【详解】} 由 $\frac{x}{x - 2} \geq  0$ ,得 $\left\{  \begin{array}{l} x\left( {x - 2}\right)  \geq  0 \\  x - 2 \neq  0 \end{array}\right.$ ,解得 $x > 2$ 或 $x \leq  0$ ,}
{\small 又 $x \in  U$ ,所以 $A = \{  - 2, - 1,0,3\}$ ,则 $\bar{A} = \{ 1,2\}$ .}
\end{solutionblock}



## 第 2 题


2. 已知一个直三棱柱的侧棱长为 2，底面面积为 2，则该三棱柱的体积为___.



\begin{solutionblock}
{\small \anslabel{【考点】：} 棱柱体积公式}
{\small \anslabel{【答案】} 4}
{\small \anslabel{【解析】}}
{\small \anslabel{【分析】} 根据三棱柱的体积公式计算求解.}
{\small \anslabel{【详解】} 设一个直三棱柱的侧棱长为 $h = 2$ ,底面面积为 $S = 2$ ,则该三棱柱的体积为 $V = {Sh} = 2 \times  2 = 4$ .}
\end{solutionblock}



## 第 3 题


3. 已知离散型随机变量 $X$ 服从二项分布 $B\left( {8,\frac{1}{2}}\right)$ ，则 $D\left( X\right)  =$ ___.



\begin{solutionblock}
{\small \anslabel{【考点】：} 二项分布的方差}
{\small \anslabel{【答案】} 2}
{\small \anslabel{【解析】}}
{\small \anslabel{【分析】} 利用二项分布的方差公式直接计算.}
{\small \anslabel{【详解】} 因为随机变量 $X$ 服从二项分布 $B\left( {8,\frac{1}{2}}\right)$ ,}
{\small 所以 $D\left( X\right)  = 8 \times  \frac{1}{2} \times  \left( {1 - \frac{1}{2}}\right)  = 2$ .}
{\small 故答案为:2.}
\end{solutionblock}



## 第 4 题


4. 设 $a : x \leq  m$ ， $\beta  : \left| {x - 2}\right|  \leq  1$ ，若 $a$ 是 $\beta$ 的必要条件，则实数 $m$ 的取值范围是___.



\begin{solutionblock}
{\small \anslabel{【考点】：} 充分必要条件、绝对值不等式}
{\small \anslabel{【答案】} $\{ m \mid  m \geq  3\}$}
{\small \anslabel{【解析】}}
{\small \anslabel{【解析】}}
{\small \anslabel{【详解】} 由 $\left| {x - 2}\right|  \leq  1$ 得 $- 1 \leq  x - 2 \leq  1$ ,得 $1 \leq  x \leq  3$ ,}
{\small 因为 $a$ 是 $\beta$ 的必要条件,所以 $\{ x \mid  1 \leq  x \leq  3\}  \subseteq  \{ x \mid  x \leq  m\}$ ,得 $m \geq  3$ ,}
{\small 故实数 $m$ 的取值范围是 $\{ m \mid  m \geq  3\}$ .}
\end{solutionblock}



## 第 5 题


5. 某工厂为判断两种不同的操作方法是否对生产某种零件的合格个数有影响，收集了相关数据，绘制了 $2 \times  2$ 列联表，设原假设 ${H}_{0}$ :两种不同的操作方法对生产该种零件的合格个数没有影响,计算出统计量 ${\chi }^{2} = {5.284}$ ,已知 $P\left( {{\chi }^{2} \geq  {3.841}}\right)  \approx  {0.05}$ ,则在显著性水平 $a = {0.05}$ 下，推断的结论为___ ${H}_{0}$ . (用“拒绝”或“接受”填空)



\begin{solutionblock}
{\small \anslabel{【考点】：} 独立性检验、假设检验}
{\small \anslabel{【答案】} 拒绝}
{\small \anslabel{【解析】}}
{\small \anslabel{【解析】}}
{\small \anslabel{【详解】} 在独立性检验中,当计算出的统计量 ${\chi }^{2}$ 大于给定显著性水平 $a$ 对应的临界值时, 样本数据出现的概率小于 $a$ ,}
{\small 属于小概率事件,根据小概率原理,我们拒绝原假设 ${H}_{0}$ ,认为两个变量之间存在显著关联, 本题中 ${\chi }^{2} = {5.284} > {3.841}$ ,所以拒绝 ${H}_{0}$ ,即认为两种操作方法对合格个数有影响.}
\end{solutionblock}



## 第 6 题


6. 在 ${ABC}$ 中， ${AB} = 1,{BC} = 2$ ， $\overrightarrow{BA} \cdot  \overrightarrow{CA} = \frac{7}{5}$ ，则 $\cos \angle {ABC} =$ ___.



\begin{solutionblock}
{\small \anslabel{【考点】：} 向量数量积、余弦定理}
{\small \anslabel{【答案】} $- \frac{1}{5}\# \#  - {0.2}$}
{\small \anslabel{【解析】}}
{\small \anslabel{【分析】} 根据数量积公式,可得 $\cos A$ 的值,根据余弦定理,可得 ${b}^{2}$ ,代入余弦定理,即可得答案.}
{\small \anslabel{【详解】} 设 ${AB} = c = 1,{BC} = a = 2,{AC} = b$ ,}
{\small 由题意 $\overrightarrow{BA} \cdot  \overrightarrow{CA} = \left| \overrightarrow{BA}\right|  \cdot  \left| \overrightarrow{CA}\right| \cos A = b\cos A = \frac{7}{5}$ ，所以 $\cos A = \frac{7}{5b}$ ，}
{\small 又 $\cos A = \frac{{b}^{2} + {c}^{2} - {a}^{2}}{2bc} = \frac{{b}^{2} + 1 - 4}{2b} = \frac{7}{5b}$ ，得 ${b}^{2} = \frac{29}{5}$ ，}
{\small 所以 $\cos \angle {ABC} = \frac{{a}^{2} + {c}^{2} - {b}^{2}}{2ac} = \frac{4 + 1 - \frac{29}{5}}{2 \times  2 \times  1} =  - \frac{1}{5}$ .}
\end{solutionblock}



## 第 7 题


7. 将 ${\left( 2 - x\right) }^{4}$ 二项展开式中的各项等可能地随机重新排列,观察排列中是否存在系数为负数的项相邻,若存在,则记随机变量 $X = 1$ ,否则记 $X = 0$ ,则 $E\left\lbrack  X\right\rbrack   =$ ___.



\begin{solutionblock}
{\small \anslabel{【考点】：} 二项式定理、排列组合、随机变量期望}
{\small \anslabel{【答案】} $\frac{2}{5}\# \# {0.4}$}
{\small \anslabel{【解析】}}
{\small \anslabel{【解析】}}
{\small \anslabel{【详解】} 二项式 ${\left( 2 - x\right) }^{4}$ 的通项公式为 ${T}_{r + 1} = {\mathrm{C}}_{4}^{r} \cdot  {2}^{4 - r} \cdot  {\left( -x\right) }^{r}$ ,}
{\small 显然该二项式展开共有 5 项，}
{\small 当 $r = 1,3$ 时，即第2,4项系数为负数，}
{\small 因为各项等可能地随机重新排列,}
{\small 所以排列数为 ${\mathrm{A}}_{5}^{5}$ ,}
{\small 系数为负数的项相邻的排列数为 ${\mathrm{A}}_{2}^{2}{\mathrm{\;A}}_{4}^{4}$ ,}
{\small 所以 $P\left( {X = 1}\right)  = \frac{{\mathrm{A}}_{2}^{2}{\mathrm{\;A}}_{4}^{4}}{{\mathrm{\;A}}_{5}^{5}} = \frac{2}{5}$ ,}
{\small 因此 $E\left\lbrack  X\right\rbrack   = 1 \times  \frac{2}{5} + 0 \times  P\left( {X = 0}\right)  = \frac{2}{5}$ .}
\end{solutionblock}



## 第 8 题


8. 在正四棱台 ${ABCD} - {A}_{1}{B}_{1}{C}_{1}{D}_{1}$ 中， ${AB} = 2,{A}_{1}{B}_{1} = 4$ ，异面直线 ${C}_{1}{D}_{1}$ 与 $A{A}_{1}$ 所成角为 $\frac{\pi }{3}$ ， 设二面角 $A - {A}_{1}{B}_{1} - {D}_{1}$ 的大小为 $\theta$ ，则 $\tan \theta  =$ ___.



\begin{solutionblock}
{\small \anslabel{【考点】：} 正四棱台、异面直线所成角、二面角}
{\small \anslabel{【答案】} $\sqrt{2}$}
{\small \anslabel{【解析】}}
{\small \anslabel{【分析】} 法一在正四棱台中,由异面直线所成角可得 $\angle A{A}_{1}{B}_{1} = {60}$ ,再根据面面角定义计算求解即可, 法二建立空间直角坐标系, 求出关键点的坐标和关键平面的法向量, 进而利用二面角的向量求法并结合同角三角函数的基本关系求解即可.}
{\small \anslabel{【详解】} 法一: 在正四棱台 ${ABCD} - {A}_{1}{B}_{1}{C}_{1}{D}_{1}$ 中, ${C}_{1}{D}_{1}//{A}_{1}{B}_{1}$ ,}
{\small 因为 ${AB} = 2$ ， ${A}_{1}{B}_{1} = 4$ ，所以 $\angle {A{A}_{1}{B}_{1}}$ 为异面直线 ${C}_{1}{D}_{1}$ 与 $A{A}_{1}$ 所成的角，}
{\small 即 $\angle A{A}_{1}{B}_{1} = {60}$ ,过点 $A$ 作平面 ${A}_{1}{B}_{1}{C}_{1}{D}_{1}$ 的垂线,垂足为 $G$ ,}
{\small 作直线 ${A}_{1}{B}_{1}$ 的垂线,垂足为 $H$ ,连接 ${GH}$ ,如图所示:}
{\small ![bo_d7obids91nqc738536r0_8_240_189_372_274_0.jpg](images/bo_d7obids91nqc738536r0_8_240_189_372_274_0.jpg)}
{\small 由正四棱台性质可知,点 $G$ 在线段 ${A}_{1}{C}_{1}$ 上, ${A}_{1}H = 1$ ,}
{\small 所以 ${AH} = \sqrt{3},{A}_{1}G = \frac{1}{2}\left( {{A}_{1}{C}_{1} - {AC}}\right)  = \sqrt{2},{GH} = \sqrt{{A}_{1}{G}^{2} - {A}_{1}{H}^{2}} = 1$ ，}
{\small 由二面角定义可知 $\angle {AHG}$ 即为二面角 $A - {A}_{1}{B}_{1} - {D}_{1}$ 的平面角,}
{\small 而 $\tan \angle {AHG} = \frac{AG}{GH} = \sqrt{2}$ ,故 $\tan \theta  = \sqrt{2}$ .}
{\small 法二: 如图,作出符合题意的图形,作下底面中心 ${O}_{1}$ ,上底面中心 $O$ , 以 ${O}_{1}$ 为原点,建立空间直角坐标系,设正四棱台的高为 $h$ ,}
{\small ![bo_d7obids91nqc738536r0_8_241_982_375_337_0.jpg](images/bo_d7obids91nqc738536r0_8_241_982_375_337_0.jpg)}
{\small 由题意得 ${AB} = 2,{A}_{1}{B}_{1} = 4$ ,则 ${C}_{1}\left( {-2,2,0}\right) ,{D}_{1}\left( {-2, - 2,0}\right) ,{B}_{1}\left( {2,2,0}\right) \; {A}_{1}\left( {2, - 2,0}\right) ,\;A\left( {1, - 1, h}\right)$ ,则 $\overrightarrow{{C}_{1}{D}_{1}} = \left( {0, - 4,0}\right) ,\;\overrightarrow{A{A}_{1}} = \left( {1, - 1, - h}\right)$ ,}
{\small 因为异面直线 ${C}_{1}{D}_{1}$ 与 $A{A}_{1}$ 所成角为 $\frac{\pi }{3}$ ,}
{\small 所以 $\frac{4}{4 \times  \sqrt{{1}^{2} + {\left( -1\right) }^{2} + {\left( -h\right) }^{2}}} = \frac{1}{2}$ ,解得 $h = \sqrt{2}$ ,}
{\small 由题意得面 ${A}_{1}{B}_{1}{D}_{1}$ 的法向量为 $\overrightarrow{n} = \left( {0,0,1}\right)$ ,}
{\small 则 $A\left( {1, - 1,\sqrt{2}}\right) ,\overrightarrow{A{A}_{1}} = \left( {1, - 1, - \sqrt{2}}\right) ,\overrightarrow{{A}_{1}{B}_{1}} = \left( {0,4,0}\right)$ ,}
{\small 设面 $A{A}_{1}{B}_{1}$ 的法向量为 $\overrightarrow{m} = \left( {x, y, z}\right)$ ,}
{\small 则 $\left\{  \begin{array}{l} \overrightarrow{A{A}_{1}} \cdot  \overrightarrow{m} = x - y - \sqrt{2}z = 0 \\  \overrightarrow{{A}_{1}{B}_{1}} \cdot  \overrightarrow{m} = {4y} = 0 \end{array}\right.$ ,令 $x = \sqrt{2}$ ,解得 $y = 0, z = 1$ ,}
{\small 得到 $\overrightarrow{m} = \left( {\sqrt{2},0,1}\right)$ ,由图可知, $\theta$ 是锐角,则 $\cos \theta  = \frac{\left| \overrightarrow{m} \cdot  \overrightarrow{n}\right| }{\left| \overrightarrow{m}\right|  \cdot  \left| \overrightarrow{n}\right| } = \frac{1}{1 \times  \sqrt{3}} = \frac{\sqrt{3}}{3}$ ,}
{\small 由已知得 $\sin \theta  > 0$ ，由同角三角函数的基本关系得 $\sin \theta  = \frac{\sqrt{6}}{3}$ ，}
{\small 故 $\tan \theta  = \frac{\sin \theta }{\cos \theta } = \frac{\frac{\sqrt{6}}{3}}{\frac{\sqrt{3}}{3}} = \sqrt{2}$ .}
\end{solutionblock}



## 第 9 题


9. 已知复数 $z$ 满足 $z + \frac{1}{z}$ 是实数，则 $\left| {z - 2 - \frac{\mathrm{i}}{2}}\right|$ 的最小值为___.



\begin{solutionblock}
{\small \anslabel{【考点】：} 复数运算、复数模的几何意义}
{\small \anslabel{【答案】} $\frac{1}{2}\# \# {0.5}$}
{\small \anslabel{【解析】}}
{\small \anslabel{【分析】} 设 $z = a + b\mathrm{i}$ ,代入 $z + \frac{1}{z}$ 中化简,由 $z + \frac{1}{z}$ 是实数,得 $b = 0$ 或 ${a}^{2} + {b}^{2} = 1$ ,利用复数模的几何意义求 $\left| {z - 2 - \frac{\mathrm{i}}{2}}\right|$ 的最小值}
{\small \anslabel{【详解】} 设 $z = a + b\mathrm{i}$ ,则 $z + \frac{1}{z} = a + b\mathrm{i} + \frac{1}{a + b\mathrm{i}} = a + b\mathrm{i} + \frac{a - b\mathrm{i}}{\left( {a + b\mathrm{i}}\right) \left( {a - b\mathrm{i}}\right) } \; = a + \frac{a}{{a}^{2} + {b}^{2}} + \left( {b - \frac{b}{{a}^{2} + {b}^{2}}}\right) \mathrm{i},$}
{\small 因为 $z + \frac{1}{z}$ 是实数,所以 $b - \frac{b}{{a}^{2} + {b}^{2}} = \frac{b\left( {{a}^{2} + {b}^{2} - 1}\right) }{{a}^{2} + {b}^{2}} = 0$ ,}
{\small 所以 $b = 0$ 或 ${a}^{2} + {b}^{2} = 1$ ,}
{\small 当 $b = 0$ 时, $z$ 的轨迹是 $x$ 轴 (除原点外),}
{\small 此时 $\left| {z - 2 - \frac{\mathrm{i}}{2}}\right|$ 的几何意义表示 $x$ 轴上的点 (除原点) 和 $\left( {2,\frac{1}{2}}\right)$ 的距离,此时 $\left| {z - 2 - \frac{\mathrm{i}}{2}}\right|  \geq  \frac{1}{2}$ ,}
{\small 当 ${a}^{2} + {b}^{2} = 1$ 时,复数 $z$ 的轨迹是以原点为圆心,1为半径的圆,如图,}
{\small ![bo_d7obids91nqc738536r0_10_238_209_337_322_0.jpg](images/bo_d7obids91nqc738536r0_10_238_209_337_322_0.jpg)}
{\small 根据复数模的几何意义可知, $\left| {z - 2 - \frac{\mathrm{i}}{2}}\right|$ 的几何意义是圆上的点到 $\left( {2,\frac{1}{2}}\right)$ 的距离,}
{\small 由图可知, $\left| {z - 2 - \frac{\mathrm{i}}{2}}\right|$ 的最小值为 $\left| {OA}\right|  - 1 = \sqrt{{2}^{2} + {\left( \frac{1}{2}\right) }^{2}} - 1 = \frac{\sqrt{17}}{2} - 1$ .}
{\small 因为 $= \frac{\sqrt{17}}{2} - 1 > \frac{1}{2}$ ,所以 $\left| {z - 2 - \frac{\mathrm{i}}{2}}\right|$ 的最小值为 $\frac{1}{2}$ .}
\end{solutionblock}



## 第 10 题


10. 已知焦点为 $F$ 的抛物线 $C : {y}^{2} = {4x}$ 上有两点 $A$ 和 $B$ ，且 $\angle {AFB} = {150}^{ \circ  }$ ， $E$ 为 $A$ 和 $B$ 的中点，过点 $E$ 作 $C$ 的准线的垂线，垂足为 $H$ ，则 $\frac{{\left| AB\right| }^{2}}{{\left| EH\right| }^{2}}$ 的最小值为___.



\begin{solutionblock}
{\small \anslabel{【考点】：} 抛物线定义、余弦定理、基本不等式}
{\small \anslabel{【答案】} $2 + \sqrt{3}$}
{\small \anslabel{【解析】}}
{\small \anslabel{【分析】} 设 $\left| {AF}\right|  = a,\left| {BF}\right|  = b$ ,根据中位线定理以及抛物线定义可得 $\left| {EH}\right|  = \frac{1}{2}\left( {a + b}\right)$ , 在 $\bigtriangleup {AFB}$ 中,由余弦定理以及基本不等式可得 $\left| {AB}\right|  \geq  \frac{2 - \sqrt{3}}{4}\left( {a + b}\right)$ ,即可求得 $\frac{{\left| AB\right| }^{2}}{{\left| EH\right| }^{2}}$ 的最小值.}
{\small \anslabel{【详解】} 设 $\left| {AF}\right|  = a,\left| {BF}\right|  = b$ ,作 ${AQ}$ 垂直抛物线的准线于点 $Q,{BP}$ 垂直抛物线的准线于点 $P$ .}
{\small ![bo_d7obids91nqc738536r0_10_240_1728_295_363_0.jpg](images/bo_d7obids91nqc738536r0_10_240_1728_295_363_0.jpg)}
{\small 由抛物线的定义,知 $\left| {AF}\right|  = \left| {AQ}\right| ,\left| {BF}\right|  = \left| {BP}\right|$ ,根据中位线定理以及抛物线定义可得 $\left| {EH}\right|  = \frac{1}{2}\left( {a + b}\right) ,$}
{\small 由余弦定理得 ${\left| AB\right| }^{2} = {a}^{2} + {b}^{2} - {2ab}\cos {150}^{ \circ  } = {a}^{2} + {b}^{2} + \sqrt{3}{ab} = {\left( a + b\right) }^{2} - \left( {2 - \sqrt{3}}\right) {ab}$ ,又 ${ab} \leq  {\left( \frac{a + b}{2}\right) }^{2}$}
{\small $\therefore {\left( a + b\right) }^{2} - \left( {2 - \sqrt{3}}\right) {ab} \geq  {\left( a + b\right) }^{2} - \frac{2 - \sqrt{3}}{4}{\left( a + b\right) }^{2} = \frac{2 + \sqrt{3}}{4}{\left( a + b\right) }^{2}$ ,当且仅当 $a = b$ 时, 等号成立,}
{\small $\therefore {\left| AB\right| }^{2} \geq  \frac{2 + \sqrt{3}}{4}{\left( a + b\right) }^{2}$ ,}
{\small $\therefore \frac{{\left| AB\right| }^{2}}{{\left| EH\right| }^{2}} \geq  \frac{\frac{2 + \sqrt{3}}{4}{\left( a + b\right) }^{2}}{{\left\lbrack  \frac{1}{2}\left( a + b\right) \right\rbrack  }^{2}} = 2 + \sqrt{3}$ ,即 $\frac{{\left| AB\right| }^{2}}{{\left| EH\right| }^{2}}$ 的最小值为 $2 + \sqrt{3}$ .}
\end{solutionblock}



## 第 11 题


11. 某种健身拉力器的手臂固定支架为 ${BE}$ ,需运动者将肩关节放置于点 $B$ 处,手肘放置于

点 $D$ 处,手掌放置于点 $E$ 处握拳握住弹力绳的一端,弹力绳的另一端连接于点 $C$ 处; 保持 $B\text{ 、 }D\text{ 、 }E\text{ 、 }C$ 四点共线. 将小臂视为线段 ${DE}$ ,长度为 $r$ ,大臂视为线段 ${BD}$ ,长度为 ${1.3r}$ . 在锻炼时要求保持肩关节 $B$ 和手肘 $D$ 不动，运用大臂力量将小臂 ${DE}$ 绕着手肘 $D$ 作圆周运动， 始终与 ${BC}$ 处于同一平面，弹力绳随之以紧绷状态从 ${CE}$ 拉伸至 ${CA}$ . 已知某位运动者健身时 ${CE} = {0.1r}$ ，当弹力绳拉至最长时， $\angle {ABC} = 3\angle {ACB}$ ，则此时 $\angle {ACB} =$ ___度.(结果精确到 0.1 度)

![bo_d7obids91nqc738536r0_1_241_879_372_197_0.jpg](images/bo_d7obids91nqc738536r0_1_241_879_372_197_0.jpg)



\begin{solutionblock}
{\small \anslabel{【考点】：} 解三角形、正弦定理、余弦定理、近似计算}
{\small \anslabel{【答案】} 16.7}
{\small \anslabel{【解析】}}
{\small \anslabel{【分析】} 先由题意得出 ${BC} = {2.4r},{CD} = {1.1r}$ ,且 ${AD} = r$ . 再设 $\angle {ACB} = x$ ,则 $\angle {ABC} = {3x}$ ,利用正弦定理表示出 ${AC}$ 的长度,最后在 $\vartriangle  {ACD}$ 中由余弦定理列方程求得 $x$ 的值.}
{\small \anslabel{【详解】} 由题意可得 $\mathrm{B},\mathrm{D},\mathrm{E},\mathrm{C}$ 四点共线, ${BD} = {1.3r},{DE} = r,{CE} = {0.1r}$ ,且点 $\mathrm{A}$ 为点 $\mathrm{E}$ 绕点 $\mathrm{D}$ 旋转后的位置,所以 ${AD} = {DE} = r$ .}
{\small 因此 ${BC} = {BD} + {DE} + {CE} = {1.3r} + r + {0.1r} = {2.4r},{CD} = {DE} + {CE} = r + {0.1r} = {1.1r}$ .}
{\small 设 $\angle {ACB} = x$ ,则 $\angle {ABC} = {3x}$ ,所以 $\angle {BAC} = {180}^{ \circ  } - {4x}$ .}
{\small 在 $\bigtriangleup  {ABC}$ 中，由正弦定理可得 $\frac{AC}{\sin {3x}} = \frac{BC}{\sin {4x}}$ ，所以 ${AC} = \frac{{2.4r}\sin {3x}}{\sin {4x}}$ .}
{\small 在 $\bigtriangleup  {ACD}$ 中，由余弦定理可得 $A{D}^{2} = A{C}^{2} + C{D}^{2} - {2AC} \cdot  {CD}\cos x$ .}
{\small 将 ${AD} = r$ ， ${CD} = {1.1r}$ ， ${AC} = \frac{{2.4r}\sin {3x}}{\sin {4x}}$ 代入，得 ${r}^{2} = {\left( \frac{{2.4r}\sin {3x}}{\sin {4x}}\right) }^{2} + {\left( {1.1}r\right) }^{2} - 2 \cdot  \frac{{2.4r}\sin {3x}}{\sin {4x}} \cdot  {1.1r}\cos x.$}
{\small 两边同除以 ${r}^{2}$ ,得 $1 = {\left( \frac{{2.4}\sin {3x}}{\sin {4x}}\right) }^{2} + {1.21} - {5.28} \cdot  \frac{\sin {3x}\cos x}{\sin {4x}}$ .}
{\small 化简可得 ${12}{\tan }^{6}x - {97}{\tan }^{4}x + {198}{\tan }^{2}x - {17} = 0$ .}
{\small 令 $t = {\tan }^{2}x$ ,则 ${12}{t}^{3} - {97}{t}^{2} + {198t} - {17} = 0$ .}
{\small 解得 $t \approx  {0.0898}$ ,或 $t \approx  {3.5588}$ ,或 $t \approx  {4.4348}$ .}
{\small 因为 $0 < x < {45}^{ \circ  }$ ,所以 $0 < {\tan }^{2}x < 1$ ,故 ${\tan }^{2}x \approx  {0.0898}$ .}
{\small 于是 $x \approx  {16.678}^{ \circ  }$ . 故此时 $\angle {ACB} \approx  {16.7}^{ \circ  }$ .}
\end{solutionblock}



## 第 12 题


12. 在以 $O$ 为原点的空间直角坐标系中,设 $\overrightarrow{i} = \left( {1,0,0}\right) ,\overrightarrow{j} = \left( {0,1,0}\right) , A$ 和 $B$ 是两个点集, 设 $A = \left\{  {P\left| {\;\overrightarrow{OP} \cdot  \overrightarrow{j} = 1}\right. ,\left\langle  {\overrightarrow{OP},\overrightarrow{j}}\right\rangle   = \frac{\pi }{4}}\right\}$ ,对任意的 $Q \in  B$ ,总存在 $P \in  A$ ,使得 $\overrightarrow{OP} \cdot  \overrightarrow{OQ} = 2$ . 若 $T \in  B,\overrightarrow{OT} = x\overrightarrow{i} + y\overrightarrow{j}\left( {x\text{ 、 }y \in  \mathbf{R}}\right)$ 且 $\left| {\overrightarrow{OT} - \overrightarrow{j}}\right|  = 2$ ,则 $\overrightarrow{OT} \cdot  \overrightarrow{i}$ 的取值范围是___.

## 二、选择题 (本大题共有 4 题, 满分 18 分, 第 13-14 题每题 4 分, 第 15-16 题 每题 5 分) 每题有且只有一个正确答案, 考生应在答题纸的相应位置上, 将所选 答案的代号涂黑>



\begin{solutionblock}
{\small \anslabel{【考点】：} 空间向量、直线与圆的位置关系}
{\small \anslabel{【答案】} $\left\lbrack  {-2,\frac{-1 - \sqrt{7}}{2}}\right\rbrack   \cup  \left\lbrack  {\frac{-1 + \sqrt{7}}{2},2}\right\rbrack$}
{\small \anslabel{【解析】}}
{\small \anslabel{【分析】} 根据空间向量数量积的坐标表示公式, 结合直线与圆的位置关系进行求解即可.}
{\small \anslabel{【详解】} 设 $P\left( {{x}_{1},{y}_{1},{z}_{1}}\right)$ ,因为 $\overrightarrow{OP} \cdot  \overrightarrow{j} = 1,\overrightarrow{OP},\overrightarrow{j} = \frac{\pi }{4}$ ,}
{\small 所以 ${y}_{1} = 1$ ,且 $\sqrt{{x}_{1}^{2} + 1 + {z}_{1}^{2}} \cdot  1 \cdot  \frac{\sqrt{2}}{2} = 1 \Rightarrow  {x}_{1}^{2} + 1 + {z}_{1}^{2} = 2 \Rightarrow  {x}_{1}^{2} + {z}_{1}^{2} = 1$ ,}
{\small 即 $P\left( {{x}_{1},1,{z}_{1}}\right)$ ,且 ${x}_{1}^{2} + {z}_{1}^{2} = 1$ ,显然 $- 1 \leq  {x}_{1} \leq  1$ ,}
{\small 设 $Q\left( {a, b, c}\right)$ ,因为 $\overrightarrow{OP} \cdot  \overrightarrow{OQ} = 2$ ,所以 $a{x}_{1} + b + c{z}_{1} = 2$ ,}
{\small 因为 $\overrightarrow{OT} = x\overrightarrow{i} + y\overrightarrow{j}\left( {x\text{ 、 }y \in  \mathbf{R}}\right)$ ,}
{\small 所以 $\overrightarrow{OT} = \left( {x, y,0}\right)$ ,}
{\small 因为 $\left| {\overrightarrow{OT} - \overrightarrow{j}}\right|  = 2$ ,}
{\small 所以 $\sqrt{{x}^{2} + {\left( y - 1\right) }^{2}} = 2 \Rightarrow  {x}^{2} + {\left( y - 1\right) }^{2} = 4$ ,}
{\small 因为 $T \in  B$ ,所以 $x{x}_{1} + y = 2 \Rightarrow  y = 2 - x{x}_{1}$ ,代入 ${x}^{2} + {\left( y - 1\right) }^{2} = 4$ 中,}
{\small 得 ${x}^{2} + {\left( 2 - x{x}_{1} - 1\right) }^{2} = 4 \Rightarrow  \left( {1 + {x}_{1}^{2}}\right) {x}^{2} - 2{x}_{1}x - 3 = 0$ ,}
{\small $\Delta  = 4{x}_{1}^{2} + {12}\left( {1 + {x}_{1}^{2}}\right)  = {12} + {16}{x}_{1}^{2} > 0,$}
{\small 因此直线 $x{x}_{1} + y - 2 = 0$ 与圆 ${x}^{2} + {\left( y - 1\right) }^{2} = 4$ 有两个不同的交点,}
{\small 因为 $- 1 \leq  {x}_{1} \leq  1$ ,所以直线 $y = 2 - x{x}_{1}$ 的斜率 $- {x}_{1}$ 的取值范围为 $- 1 \leq   - {x}_{1} \leq  1$ , 如下图所示:}
{\small ![bo_d7obids91nqc738536r0_13_239_1348_349_342_0.jpg](images/bo_d7obids91nqc738536r0_13_239_1348_349_342_0.jpg)}
{\small 由 $\left\{  {\begin{array}{l} {x}^{2} + {\left( y - 1\right) }^{2} = 4 \\  y = x + 2 \end{array} \Rightarrow  2{x}^{2} + {2x} - 3 = 0 \Rightarrow  x = \frac{-1 \pm  \sqrt{7}}{2}}\right.$}
{\small 直线 $y = x + 2$ 与圆 ${x}^{2} + {\left( y - 1\right) }^{2} = 4$ 的交点坐标为 $A\left( {\frac{-1 - \sqrt{7}}{2},0}\right) , B\left( {\frac{-1 + \sqrt{7}}{2},0}\right)$ , 又因为直线 $y = x + 2$ 、直线 $y =  - x + 2$ 斜率互为相反数,且过同一点 $\left( {0,2}\right)$ ,与圆 ${x}^{2} + {\left( y - 1\right) }^{2} = 4$ 都是关于纵轴对称,}
{\small 所以当 $- 1 \leq   - {x}_{1} \leq  1$ 时, $x \in  \left\lbrack  {-2,\frac{-1 - \sqrt{7}}{2}}\right\rbrack   \cup  \left\lbrack  {\frac{-1 + \sqrt{7}}{2},2}\right\rbrack$}
{\small 因为 $\overrightarrow{OT} \cdot  \overrightarrow{i} = x$ ,}
{\small 所以 $\overrightarrow{OT} \cdot  \overrightarrow{i}$ 的取值范围是 $\left\lbrack  {-2,\frac{-1 - \sqrt{7}}{2}}\right\rbrack   \cup  \left\lbrack  {\frac{-1 + \sqrt{7}}{2},2}\right\rbrack$ .}
\end{solutionblock}



## 第 13 题


13. 双曲线 ${y}^{2} - 2{x}^{2} = 1$ 的渐近线是( ).

A. $y =  \pm  \sqrt{2}x$ B. $x =  \pm  \sqrt{2}y$ C. $y =  \pm  {2x}$ D. $x =  \pm  {2y}$



\begin{solutionblock}
{\small \anslabel{【考点】：} 双曲线的渐近线}
{\small \anslabel{【答案】} A}
{\small \anslabel{【解析】}}
{\small \anslabel{【解析】}}
{\small \anslabel{【详解】} 双曲线 ${y}^{2} - 2{x}^{2} = 1$ 的标准形式为 ${y}^{2} - \frac{{x}^{2}}{\frac{1}{2}} = 1$ ,}
{\small 显然该双曲线焦点在 $y$ 轴上,其中 ${a}^{2} = 1,{b}^{2} = \frac{1}{2}$ ,即 $a = 1, b = \frac{\sqrt{2}}{2}$ ,}
{\small 因为焦点在 $y$ 轴上的双曲线的渐近线公式是 $y =  \pm  \frac{a}{b}x$ ,且 $\frac{a}{b} = \frac{1}{\frac{\sqrt{2}}{2}} = \sqrt{2}$ ,}
{\small 所以双曲线 ${y}^{2} - 2{x}^{2} = 1$ 的渐近线是 $y =  \pm  \sqrt{2}x$ .}
\end{solutionblock}



## 第 14 题


14. 已知 $P\left( A\right)$ 和 $P\left( B\right)$ 分别表示事件 $A$ 和事件 $B$ 发生的概率,且 $0 < P\left( B\right)  < 1$ ,则在下列各项中，“ $A$ 和 $B$ 独立” 的充分条件是( ).

A. $P\left( A\right)  = 1 - P\left( B\right)$ B. $P\left( {A \mid  B}\right)  = P\left( {A \mid  \bar{B}}\right)$

C. $P\left( {A \cup  B}\right)  = P\left( A\right)  + P\left( B\right)$ D. $P\left( {A \mid  B}\right)  = P\left( {\bar{A} \mid  B}\right)$



\begin{solutionblock}
{\small \anslabel{【考点】：} 事件独立性、条件概率}
{\small \anslabel{【答案】} B}
{\small \anslabel{【解析】}}
{\small \anslabel{【分析】} 事件 $A$ 和 $B$ 独立的定义是 $P\left( {AB}\right)  = P\left( A\right) P\left( B\right)$ ,根据每个选项的条件,结合条件概率公式 $P\left( {A \mid  B}\right)  = \frac{P\left( {AB}\right) }{P\left( B\right) }$ 以及概率的基本性质,判断是否能推出 $P\left( {AB}\right)  = P\left( A\right) P\left( B\right)$ .}
{\small \anslabel{【详解】} 选项 $\mathrm{A}, P\left( A\right)  = 1 - P\left( B\right)$ ,则 $P\left( A\right)  + P\left( B\right)  = 1$ ,并不能推出 $P\left( {AB}\right)  = P\left( A\right) P\left( B\right)$ , 所以事件 $A$ 和 $B$ 不一定独立, $\mathrm{A}$ 错误.}
{\small 选项 B, $P\left( {A \mid  B}\right)  = \frac{P\left( {AB}\right) }{P\left( B\right) }, P\left( {A \mid  \bar{B}}\right)  = \frac{P\left( {A\bar{B}}\right) }{P\left( \bar{B}\right) },\because P\left( {A \mid  B}\right)  = P\left( {A \mid  \bar{B}}\right)$ , $\therefore \frac{P\left( {AB}\right) }{P\left( B\right) } = \frac{P\left( {A\bar{B}}\right) }{P\left( \bar{B}\right) }$ .}
{\small $\because P\left( {A\bar{B}}\right)  = P\left( A\right)  - P\left( {AB}\right) , P\left( \bar{B}\right)  = 1 - P\left( B\right) ,\;\therefore \frac{P\left( {AB}\right) }{P\left( B\right) } = \frac{P\left( A\right)  - P\left( {AB}\right) }{1 - P\left( B\right) }$ .}
{\small $\therefore P\left( {AB}\right)  - P\left( {AB}\right) P\left( B\right)  = P\left( B\right) P\left( A\right)  - P\left( B\right) P\left( {AB}\right)$ .}
{\small $\therefore P\left( {AB}\right)  = P\left( A\right) P\left( B\right)$ ,即 $A$ 和 $B$ 独立, $\mathrm{B}$ 选项正确.}
{\small 选项 C, $P\left( {A \cup  B}\right)  = P\left( A\right)  + P\left( B\right)  - P\left( {AB}\right)$ ,又 $P\left( {A \cup  B}\right)  = P\left( A\right)  + P\left( B\right) , P\left( {AB}\right)  = 0$ ,}
{\small $\therefore A$ 和 $B$ 不一定独立, $\mathrm{C}$ 错误.}
{\small 选项 D, $P\left( {A \mid  B}\right)  = \frac{P\left( {AB}\right) }{P\left( B\right) }, P\left( {\bar{A} \mid  B}\right)  = \frac{P\left( {\bar{A}B}\right) }{P\left( B\right) }, P\left( {A \mid  B}\right)  = P\left( {\bar{A} \mid  B}\right) ,\therefore P\left( {AB}\right)  = P\left( {\bar{A}B}\right)$ . 又 $P\left( {\bar{A}B}\right)  = P\left( B\right)  - P\left( {AB}\right) ,\therefore P\left( {AB}\right)  = P\left( B\right)  - P\left( {AB}\right)$ ,可得 $P\left( {AB}\right)  = \frac{1}{2}P\left( B\right)$ ,也不能推出 $A$ 和 $B$ 独立, $\mathrm{D}$ 错误.}
\end{solutionblock}



## 第 15 题


15. 一定存在各项均为正数且不为常数列的无穷等差数列 $\left\{  {a}_{n}\right\}$ ,使得 ( ).

A. $\left\{  {\sin {a}_{n}}\right\}$ 为严格增数列 B. $\left\{  {\cos {a}_{n}}\right\}$ 为公差不为零的等差数

列.

C. $\left\{  {\cos {S}_{n}}\right\}$ 为等比数列 (其中, ${S}_{n} = \mathop{\sum }\limits_{{i = 1}}^{n}{a}_{i}$ ) D. $\left\{  {\sin \frac{1}{{a}_{n}}}\right\}$ 为周期数列



\begin{solutionblock}
{\small \anslabel{【考点】：} 等差数列、三角函数性质、数列综合}
{\small \anslabel{【答案】} C}
{\small \anslabel{【解析】}}
{\small \anslabel{【分析】} 利用等差数列、等比数列及三角函数性质对各选项逐一分析是否存在符合条件的无穷等差数列.}
{\small \anslabel{【详解】} 在 $\mathrm{A}$ 选项中，假设 $\left\{  {\sin {a}_{n}}\right\}$ 是严格递增的无穷数列,}
{\small 但 $\sin {a}_{n} \in  \left\lbrack  {-1,1}\right\rbrack$ ,且是周期函数,在一个周期内有增有减,}
{\small 对无穷等差数列 $\left\{  {a}_{n}\right\}$ ,当 $n \rightarrow   + \infty$ ,则 ${a}_{n} \rightarrow   + \infty$ ,}
{\small 所以 $\sin {a}_{n}$ 会周期性波动,不可能一直严格递增, A 错误,}
{\small 在 $\mathrm{B}$ 选项中，设等差数列 $\left\{  {a}_{n}\right\}$ 的首项为 ${a}_{1}$ ，公差为 $d\left( {d \neq  0}\right)$ ，}
{\small 则 ${a}_{n} = {a}_{1} + \left( {n - 1}\right) d$ ,所以}
{\small $\cos {a}_{n + 1} - \cos {a}_{n} =  - 2\sin \frac{{a}_{n + 1} + {a}_{n}}{2}\sin \frac{{a}_{n + 1} - {a}_{n}}{2} =  - 2\sin \frac{2{a}_{1} + \left( {{2n} - 1}\right) d}{2}\sin \frac{d}{2}$ ,}
{\small 由于 $\sin \frac{2{a}_{1} + \left( {{2n} - 1}\right) d}{2}$ 是关于 $n$ 的周期变化的函数,}
{\small 所以 $\cos {a}_{n + 1} - \cos {a}_{n}$ 不是常数,}
{\small 即 $\left\{  {\cos {a}_{n}}\right\}$ 不是公差不为零的等差数列, B 错误,}
{\small 在 $\mathrm{C}$ 选项中,设等差数列 $\left\{  {a}_{n}\right\}$ 的首项为 ${a}_{1}$ ,公差为 $d\left( {d \neq  0}\right)$ ,}
{\small 则 ${S}_{n} = n{a}_{1} + \frac{n\left( {n - 1}\right) d}{2}$ ,}
{\small 若 $\left\{  {\cos {S}_{n}}\right\}$ 为等比数列,则 ${\left( \cos {S}_{2}\right) }^{2} = \cos {S}_{1}\cos {S}_{3}$ ,}
{\small ${S}_{1} = {a}_{1},\;{S}_{2} = 2{a}_{1} + \frac{2 \times  1}{2}d = 2{a}_{1} + d,\;{S}_{3} = 3{a}_{1} + \frac{3 \times  2}{2}d = 3{a}_{1} + {3d},$}
{\small ${\cos }^{2}\left( {2{a}_{1} + d}\right)  = \cos {a}_{1}\cos \left( {3{a}_{1} + {3d}}\right) ,$}
{\small 当 ${a}_{1} = \pi , d = {2\pi }$ 时,则 ${a}_{n} = \pi  + \left( {\mathrm{n} - 1}\right) {2\pi } = \left( {2\mathrm{n} - 1}\right) \pi$ ,}
{\small 该数列各项均为正数，且不为常数列，其 $n$ 项和为}
{\small ${S}_{n} = \frac{n\left( {{a}_{1} + {a}_{n}}\right) }{2} = \frac{n\left\lbrack  {\pi  + \left( {2\mathrm{n} - 1}\right) \pi }\right\rbrack  }{2} = {n}^{2}\pi$ ,此时数列 $\cos {S}_{n} = \cos \left( {{n}^{2}\pi }\right)$ ,}
{\small 当 $n$ 为奇数时, ${n}^{2}$ 为奇数, $\cos \left( {{n}^{2}\pi }\right)  =  - 1$ ,}
{\small 当 $n$ 为偶数时, ${n}^{2}$ 为偶数, $\cos \left( {{n}^{2}\pi }\right)  = 1$ ,}
{\small 故 $\left\{  {\cos {S}_{n}}\right\}$ 为数列 $\{  - 1,1, - 1,1,\cdots \}$ ，是公比为 -1 的等比数列，}
{\small 所以存在这样的无穷等差数列 $\left\{  {a}_{n}\right\}$ 使得 $\left\{  {\cos {S}_{n}}\right\}$ 为等比数列, C 正确,}
{\small 在 $\mathrm{D}$ 选项中，因为 $\left\{  {a}_{n}\right\}$ 是各项为正数且不为常数列的无穷等差数列,}
{\small 所以 $\mathop{\lim }\limits_{{n \rightarrow  \infty }}{a}_{n} =  + \infty$ ,即 $\mathop{\lim }\limits_{{n \rightarrow  \infty }}\frac{1}{{a}_{n}} = 0$ ,}
{\small 所以 $\sin \frac{1}{{a}_{n}}$ 会趋近于 $\sin 0 = 0$ ,}
{\small 而周期数列是指经过一定的项数后会重复出现相同的项,}
{\small 所以 $\left\{  {\sin \frac{1}{{a}_{n}}}\right\}$ 不可能是周期数列, D 错误.}
\end{solutionblock}



## 第 16 题


16. 设 $y = f\left( x\right)$ 和 $y = g\left( x\right)$ 是两个不同的函数,且定义域和值域均为 $\mathbf{R}$ ,设 $M = \left\{  {\left( {a, b}\right)  \mid  f\left( a\right) \rangle g\left( b\right) , a, b \in  \mathbf{R}}\right\}$ ,则对于以下两个结论,说法正确的是 ( )

结论①:若当 $\left( {a, b}\right)  \in  M$ ，恒有 $\left( {-a, b}\right)  \in  M$ ，则函数 $y = f\left( x\right)$ 一定是偶函数；

结论②:若当 $\left( {a, b}\right)  \in  M$ ，恒有 $\left( {-a, - b}\right)  \in  M$ ，则函数 $y = g\left( x\right)$ 可以不是偶函数.

A. ①和②都正确 B. ①正确，②错误 C. ①错误，②正确 D. ①和② 都错误

## 三、解答题



\begin{solutionblock}
{\small \anslabel{【考点】：} 函数性质、集合、逻辑推理}
{\small \anslabel{【答案】} B}
{\small \anslabel{【解析】}}
{\small \anslabel{【分析】} 对于结论①,利用反证法假设存在 $f\left( a\right)  > f\left( {-a}\right)$ ,找到 $g\left( b\right)$ 满足 $f\left( a\right)  > g\left( b\right)  > f\left( {-a}\right)$ ,引出矛盾即可证明正确; 对于结论②,采用类似的分析找到 $f\left( a\right)$ 满足 $g\left( b\right)  > f\left( a\right)  > g\left( {-b}\right)$ ,利用 $f\left( a\right)  > g\left( {-b}\right)$ 推出 $f\left( {-a}\right)  > g\left( b\right)  > g\left( {-b}\right)$ ,再利用 $f\left( {-a}\right)  > g\left( {-b}\right)$ 推出 $f\left( a\right)  > g\left( b\right)$ ,引出矛盾即可证明错误.}
{\small \anslabel{【详解】} 对于结论①,若函数 $y = f\left( x\right)$ 不是偶函数,则存在 $f\left( a\right)  \neq  f\left( {-a}\right)$ ,}
{\small 不妨设 $f\left( a\right)  > f\left( {-a}\right)$ (否则用 $- a$ 取代 $a$ ),因为 $y = f\left( x\right)$ 和 $y = g\left( x\right)$ 值域均为 $\mathbf{R}$ ,}
{\small 则存在 $b \in  \mathbf{R}$ 使得 $g\left( b\right)  = \frac{f\left( a\right)  + f\left( {-a}\right) }{2}$ ,此时有 $f\left( a\right)  > g\left( b\right)  > f\left( {-a}\right)$ ,}
{\small 根据 $f\left( a\right)  > g\left( b\right)$ ,依题意有 $f\left( {-a}\right)  > g\left( b\right)$ ,这与 $g\left( b\right)  > f\left( {-a}\right)$ 矛盾,}
{\small 故函数 $y = f\left( x\right)$ 一定是偶函数，结论①正确；}
{\small 对于结论②，若函数 $y = g\left( x\right)$ 不是偶函数，则存在 $g\left( b\right)  \neq  g\left( {-b}\right)$ ，}
{\small 不妨设 $g\left( b\right)  > g\left( {-b}\right)$ (否则用 $- b$ 取代 $b$ ),因为 $y = f\left( x\right)$ 和 $y = g\left( x\right)$ 值域均为 $\mathbf{R}$ ,}
{\small 则存在 $a \in  \mathbf{R}$ 使得 $f\left( a\right)  = \frac{g\left( b\right)  + g\left( {-b}\right) }{2}$ ,此时 $g\left( b\right)  > f\left( a\right)  > g\left( {-b}\right)$ ,}
{\small 依题意,由 $f\left( a\right)  > g\left( {-b}\right)$ 有 $f\left( {-a}\right)  > g\left( {-\left( {-b}\right) }\right)$ ,即 $f\left( {-a}\right)  > g\left( b\right)$ ,所以}
{\small $f\left( {-a}\right)  > g\left( {-b}\right) ,$}
{\small 而 $f\left( {-a}\right)  > g\left( {-b}\right)$ 可推出 $f\left( {-\left( {-a}\right) }\right)  > g\left( {-\left( {-b}\right) }\right)$ 即 $f\left( a\right)  > g\left( b\right)$ ,与 $g\left( b\right)  > f\left( a\right)$ 矛盾,}
{\small 故函数 $y = g\left( x\right)$ 一定是偶函数,结论②错误.}
{\small \anslabel{【点睛】} 本题采用了反证法证明奇偶性, 通过灵活利用已知条件得到矛盾的结果, 并利用了两个实数之间总能找到一个实数这一结论.}
\end{solutionblock}



## 第 17 题


17. 如图,在底面半径为 2,侧面积为 $4\sqrt{5}\pi$ 的圆锥

![bo_d7obids91nqc738536r0_2_244_1200_249_334_0.jpg](images/bo_d7obids91nqc738536r0_2_244_1200_249_334_0.jpg)

$\mathrm{P} - \mathrm{O}$ 中, $A\text{ 、 }B\text{ 、 }C$ 为底面圆周上不同的三个点, $\angle {AOB} = \angle {BOC} = {45}^{ \circ  }$ .

(1)求直线 ${OB}$ 与平面 ${PAC}$ 所成角的正弦值

(2)设点 $D$ 为线段 ${PB}$ 上的动点(不含端点 $P$ 和 $B$ )，求证直线 ${OA}$ 与 ${CD}$ 不垂直



\begin{solutionblock}
{\small \anslabel{【考点】：} 圆锥的侧面积、空间向量、线面角、线线垂直}
{\small \anslabel{【答案】} (1) $\frac{2\sqrt{2}}{3}$}
{\small (2)证明见解析}
{\small \anslabel{【解析】}}
{\small \anslabel{【分析】} (1) 利用侧面积求解出圆锥的高,再利用建系的方法,求出直线 ${OB}$ 与平面 ${PAC}$ 所成角的正弦值.}
{\small (2)利用向量的方法，将点 $D$ 在线段 ${PB}$ 上刻画为 $\overrightarrow{PD} = \lambda \overrightarrow{PB}$ . 接着证明直线 $\overrightarrow{OA}$ 与 $\overrightarrow{CD}$ 的数量积不为零.}
{\small \textbf{【小问 1 详解】}}
{\small ![bo_d7obids91nqc738536r0_19_242_812_292_419_0.jpg](images/bo_d7obids91nqc738536r0_19_242_812_292_419_0.jpg)}
{\small $\because \angle {AOB} = \angle {BOC} = {45}^{ \circ  }$ ,}
{\small $\therefore \angle {AOC} = {90}^{ \circ  }$ ,}
{\small 所以可以建立如右图所示的空间直角坐标系,}
{\small 所以 $B\left( {\sqrt{2},\sqrt{2},0}\right) ,\overrightarrow{OB} = \left( {\sqrt{2},\sqrt{2},0}\right)$ ,设 ${OP} = a$ ,}
{\small 则母线长为 $\sqrt{4 + {a}^{2}}$}
{\small 则侧面积为 $\frac{1}{2}\sqrt{4 + {a}^{2}} \times  {2\pi } \times  2 = 2\sqrt{4 + {a}^{2}}\pi  = 4\sqrt{5}\pi$ ,}
{\small 解得: $a = 4$ ，}
{\small 所以 $P\left( {0,0,4}\right)$ ,}
{\small $A\left( {2,0,0}\right) , C\left( {0,2,0}\right) ,\overrightarrow{AC} = \left( {-2,2,0}\right) ,\overrightarrow{PA} = \left( {2,0, - 4}\right)$}
{\small 设面 ${PAC}$ 的法向量为 $\overrightarrow{m} = \left( {{x}_{1},{y}_{1},{z}_{1}}\right)$}
{\small 所以 $\overrightarrow{m} \cdot  \overrightarrow{AC} =  - 2{x}_{1} + 2{y}_{1} = 0,\overrightarrow{m} \cdot  \overrightarrow{PA} = 2{x}_{1} - 4{z}_{1} = 0$ ,}
{\small 不妨令 ${x}_{1} = {y}_{1} = 1$ ,则 ${z}_{1} = \frac{1}{2}$ ,}
{\small 所以 $\overrightarrow{m} = \left( {1,1,\frac{1}{2}}\right)$ ,设直线 ${OB}$ 与平面 ${PAC}$ 所成角为 $\theta$ ,}
{\small 所以 $\sin \theta  = \left| {\cos \langle \overrightarrow{OB},\overrightarrow{m}\rangle }\right|  = \frac{\left| \overrightarrow{OB} \cdot  \overrightarrow{m}\right| }{\left| \overrightarrow{OB}\right|  \cdot  \left| \overrightarrow{m}\right| } = \frac{2\sqrt{2}}{2 \times  \sqrt{{1}^{2} + {1}^{2} + {\left( \frac{1}{2}\right) }^{2}}} = \frac{2\sqrt{2}}{3}$ .}
{\small \textbf{【小问 2 详解】}}
{\small 因为 $D$ 为线段 ${PB}$ 上的动点,所以设 $\overrightarrow{PD} = \lambda \overrightarrow{PB},\lambda  \in  \left( {0,1}\right)$ ,}
{\small $\overrightarrow{CD} = \overrightarrow{CP} + \overrightarrow{PD}$ ,由 (1) 知 $\overrightarrow{PD} = \lambda \overrightarrow{PB} = \lambda \left( {\sqrt{2},\sqrt{2}, - 4}\right)  = \left( {\sqrt{2}\lambda ,\sqrt{2}\lambda , - {4\lambda }}\right)$ ,}
{\small $\overrightarrow{CP} = \left( {0, - 2,4}\right) ,\overrightarrow{CD} = \overrightarrow{CP} + \overrightarrow{PD} = \left( {\sqrt{2}\lambda ,\sqrt{2}\lambda  - 2,4 - {4\lambda }}\right) ,$}
{\small $\overrightarrow{OA} = \left( {2,0,0}\right) ,$}
{\small $\overrightarrow{OA} \cdot  \overrightarrow{CD} = 2\sqrt{2}\lambda  \neq  0,$}
{\small 所以直线 ${OA}$ 与 ${CD}$ 不垂直.}
\end{solutionblock}



## 第 18 题


18. 班主任小明查阅了某大学发表的一项本市高三学生手机使用情况的研究报告. 报告指出, 高三学生每周手机使用时长 $X$ (单位:小时) 总体上服从正态分布 $N\left( {{12},{4}^{2}}\right)$ .

(1)小明老师将自己所带班级(共 50 名学生)视为从本市高三学生总体中随机抽取的一个样本，能以此正态分布模型估算出全班每周平均手机使用时长超过 16 小时的人数，在此估

算基础上若在全班任选 3 位同学，则至少有 2 位同学的每周手机使用时长超过 16 小时的概率是多少? (结果用最简分数表示)

参考数据: 若 $X \sim  N\left( {\mu ,{\sigma }^{2}}\right)$ ,则 $P\left( {\left| {X - \mu }\right|  \leq  \sigma }\right)  \approx  {68}\%$ .

(2)小明老师发现小虹同学每周手机使用时长超过 16 小时，对其进行疏导劝解，并跟进统计出之后 5 周小虹每周手机使用时长 $x$ 与该周数学练习得分 $y$ (每周练习的难度相同且满分均为 150 分),制成表 1.以这 5 组数据建立回归方程 $y =  - {3.3x} + {146}$ . 请求出实数 $m$ 的值表 1

\begin{tabular}{|c|c|c|c|c|c|}
\hline
 & 第 1 周 & 第 2 周 & 第 3 周 & 第 4 周 & 第 5 周 \\ \hline
手机使用时长 $x$ & 20 & 18 & 22 & 16 & 14 \\ \hline
练习得分 $y$ & 80 & 88 & 73 & 92 & $m$ \\ \hline
\end{tabular}

(3)受到鼓励的小虹制定了寒假复习计划表递交给小明老师，严格控制手机使用时长.小明老师统计发现该计划表中若第 $n$ 天能复习时长超过 5 小时 (记为“高效复习”),则第 $n + 1$ 天也能“高效复习”的概率为 $\frac{3}{4}$ ; 若第 $n$ 天不能“高效复习”,则第 $n + 1$ 天还能“高效复习”的概率为 $\frac{5}{7}$ . 设 ${P}_{n}\left( {1 \leq  n \leq  {28}}\right.$ , $n$ 为正整数) 表示第 $n$ 天能 “高效复习”的概率, ${P}_{1} = 1$ ,若 ${P}_{n} > {0.7}$ 表示复习计划表第 $n$ 天有效. 求证: 数列 $\left\{  {{P}_{n} - \frac{20}{27}}\right\}$ 是等比数列,并说明小虹的该复习计划表是否在寒假每一天均有效.



\begin{solutionblock}
{\small \anslabel{【考点】：} 正态分布、古典概型、线性回归、数列递推}
{\small \anslabel{【答案】} (1) $\frac{11}{175}$}
{\small (2) 100 (3)答案见解析}
{\small \anslabel{【解析】}}
{\small \anslabel{【分析】} (1) 根据正态分布的性质和概率相关知识计算即可.}
{\small (2)先求出 $x, y$ 的平均值，然后代入回归方程即可求出结果.}
{\small (3)先根据题意列出递推式,然后证明数列 $\left\{  {{P}_{n} - \frac{20}{27}}\right\}$ 是以 $\frac{1}{28}$ 为公比的等比数列,进而可根据等比数列的通项公式求出 ${P}_{n}$ ,并根据 $n$ 的范围证明结论即可.}
{\small \textbf{【小问 1 详解】}}
{\small 由题意知,因为 $P\left( {\left| {X - {12}}\right|  \leq  4}\right)  = P\left( {8 \leq  X \leq  {16}}\right)  = {68}\%$ .}
{\small 所以任取 1 人使用手机超过 16 小时的概率为 $P = \frac{1 - {68}\% }{2} = {0.16}$ ,}
{\small 50 名同学中有 ${50} \times  {0.16} = 8$ 位超过 16 小时,}
{\small 那么至少 2 位同学使用手机超过 16 小时的概率为 $P = \frac{{\mathrm{C}}_{8}^{2} \cdot  {\mathrm{C}}_{42}^{1} + {\mathrm{C}}_{8}^{3}}{{\mathrm{C}}_{50}^{3}} = \frac{11}{175}$ .}
{\small \textbf{【小问 2 详解】}}
{\small 由题意得 $\bar{x} = \frac{1}{5}\left( {{20} + {18} + {22} + {16} + {14}}\right)  = {18},\bar{y} = \frac{1}{5}\left( {{80} + {88} + {73} + {92} + m}\right)  = \frac{{333} + m}{5}$ . 代入回归方程有 $\frac{{333} + m}{5} =  - {3.3} \times  {18} + {146}$ ,解得 $m = {100}$ .}
{\small \textbf{【小问 3 详解】}}
{\small 证明: 由题意知 ${P}_{n + 1} = \frac{3}{4}{P}_{n} + \frac{5}{7}\left( {1 - {P}_{n}}\right)  = \frac{1}{28}{P}_{n} + \frac{5}{7}$ ,}
{\small 所以 ${P}_{n + 1} - \frac{20}{27} = \frac{1}{28}{P}_{n} + \frac{5}{7} - \frac{20}{27} = \frac{1}{28}{P}_{n} - \frac{5}{189} = \frac{1}{28}\left( {{P}_{n} - \frac{20}{27}}\right)$}
{\small 所以 $\left\{  {{P}_{n} - \frac{20}{27}}\right\}$ 是以 $\frac{1}{28}$ 为公比的等比数列.}
{\small 所以 ${P}_{n} - \frac{20}{27} = \left( {{P}_{1} - \frac{20}{27}}\right) {\left( \frac{1}{28}\right) }^{n - 1} \Rightarrow  {P}_{n} = \frac{7}{27}{\left( \frac{1}{28}\right) }^{n - 1} + \frac{20}{27}$ .}
{\small 因为 $n \in  \left\lbrack  {1,{28}}\right\rbrack$ 时, ${\left( \frac{1}{28}\right) }^{n - 1} > 0$ 恒成立,所以 ${P}_{n} > \frac{20}{27} > {0.7}$ .}
{\small 所以小虹的该复习计划表在寒假每一天均有效.}
\end{solutionblock}



## 第 19 题


19. 设 $f\left( x\right)  = {\mathrm{e}}^{x} - {\mathrm{e}}^{-x}$ .

(1)解不等式: $f\left( {\ln x}\right)  < 2$ ；

(2)设 $g\left( x\right)  = f\left( x\right)  + \sin {2x}$ ，若存在 $x \in  \left\lbrack  {-2,2}\right\rbrack$ ，使得 $g\left( x\right)  + g\left( {{x}^{2} + a}\right)  > 0$ ，求实数 $a$ 的取值范围.



\begin{solutionblock}
{\small \anslabel{【考点】：} 指数函数、不等式求解、函数奇偶性与单调性}
{\small \anslabel{【答案】} (1) $\left\{  {x\left| {\;0 < x < \sqrt{2} + 1}\right. }\right\}$}
{\small (2) $\left( {-6, + \infty }\right)$}
{\small \anslabel{【解析】}}
{\small \anslabel{【分析】} (1) 化简 $f\left( {\ln x}\right)  < 2$ ,解一元二次不等式即可;}
{\small (2)先求证 $g\left( x\right)$ 的奇偶性和单调性,将问题转化为存在 $x \in  \left\lbrack  {-2,2}\right\rbrack$ ,使得 $a >  - {x}^{2} - x$ ,求一元二次函数的最小值即可.}
{\small \textbf{【小问 1 详解】}}
{\small 因为 $f\left( {\ln x}\right)  = {\mathrm{e}}^{\ln x} - {\mathrm{e}}^{-\ln x} = x - \frac{1}{x} < 2, x > 0$ ,}
{\small 所以 ${x}^{2} - {2x} - 1 < 0$ ,得 $0 < x < \sqrt{2} + 1$ , 故 $f\left( {\ln x}\right)  < 2$ 的解集为 $\left\{  {x\left| {\;0 < x < \sqrt{2} + 1}\right. }\right\}$ ;}
{\small \textbf{【小问 2 详解】}}
{\small $g\left( x\right)  = f\left( x\right)  + \sin {2x} = {\mathrm{e}}^{x} - {\mathrm{e}}^{-x} + \sin {2x}$ ,则 $g\left( {-x}\right)  = {\mathrm{e}}^{-x} - {\mathrm{e}}^{x} - \sin {2x} =  - g\left( x\right)$ ,}
{\small 因为 $x \in  \left\lbrack  {-2,2}\right\rbrack$ 关于原点对称,所以 $g\left( x\right)$ 在 $x \in  \left\lbrack  {-2,2}\right\rbrack$ 上为奇函数,}
{\small 易得 ${g}^{\prime }\left( x\right)  = {\mathrm{e}}^{x} + {\mathrm{e}}^{-x} + 2\cos {2x}$ ,}
{\small 因为 ${\mathrm{e}}^{x} + {\mathrm{e}}^{-x} \geq  2\sqrt{{\mathrm{e}}^{x} \cdot  {\mathrm{e}}^{-x}} = 2$ ,等号成立时 $x = 0$ ,所以 ${g}^{\prime }\left( x\right)  \geq  2\left( {1 + \cos {2x}}\right)  \geq  0$ , 则 $g\left( x\right)$ 在 $x \in  \left\lbrack  {-2,2}\right\rbrack$ 上单调递增,}
{\small 若存在 $x \in  \left\lbrack  {-2,2}\right\rbrack$ ,使得 $g\left( x\right)  + g\left( {{x}^{2} + a}\right)  > 0$ ,}
{\small 则存在 $x \in  \left\lbrack  {-2,2}\right\rbrack$ ,使得 $g\left( {{x}^{2} + a}\right)  >  - g\left( x\right)  = g\left( {-x}\right)$ ,}
{\small 则存在 $x \in  \left\lbrack  {-2,2}\right\rbrack$ ,使得 ${x}^{2} + a >  - x$ ,即 $a >  - {x}^{2} - x$ ,}
{\small 因为 $y =  - {x}^{2} - x$ 函数图象关于 $x =  - \frac{1}{2}$ 对称,其在 $\left\lbrack  {-2,2}\right\rbrack$ 上的最小值为 ${y}_{\min } =  - 6$ ,则 $a >  - 6$ ,}
{\small 故实数 $a$ 的取值范围为 $\left( {-6, + \infty }\right)$ .}
\end{solutionblock}



## 第 20 题


20. 设椭圆 $\Gamma  : \frac{{x}^{2}}{4} + {y}^{2} = 1$ 的左顶点为 $A$ .

(1)求 $\Gamma$ 的离心率；

(2)设 $\Gamma$ 的左焦点为 $F$ ，上顶点为 $B$ ，若点 $P$ 在 $\Gamma$ 上且位于 $y$ 轴右侧. $\overrightarrow{AB}//\overrightarrow{FP}$ ，求点 $P$ 的横坐标;

(3)设直线 $l : x - y + m = 0$ ， $l$ 与 $\Gamma$ 交于不同的两点 $C$ 和 $D$ ，若点 $A$ 在以 ${CD}$ 为直径的圆外,求实数 $m$ 的取值范围.



\begin{solutionblock}
{\small \anslabel{【考点】：} 椭圆离心率、向量平行、直线与椭圆位置关系、圆}
{\small \anslabel{【答案】} (1) $\frac{\sqrt{3}}{2}$}
{\small (2) $x = \frac{\sqrt{5} - \sqrt{3}}{2}$}
{\small (3) $m \in  \left( {-\sqrt{5},\frac{6}{5}}\right)  \cup  \left( {2,\sqrt{5}}\right)$}
{\small \anslabel{【解析】}}
{\small \anslabel{【解析】}}
{\small \textbf{【小问 1 详解】}}
{\small 由椭圆方程可得 $a = 2, b = 1, c = \sqrt{{a}^{2} - {b}^{2}} = \sqrt{3}$ ,所以 $e = \frac{c}{a} = \frac{\sqrt{3}}{2}$ .}
{\small \textbf{【小问 2 详解】}}
{\small 由条件可知 $A\left( {-2,0}\right) , B\left( {0,1}\right) , F\left( {-\sqrt{3},0}\right)$ ,}
{\small 设直线 ${AB}$ 的斜率为 ${k}_{AB}$ ,直线 ${FP}$ 的斜率为 ${k}_{FP}$ ,}
{\small ${k}_{AB} = \frac{1 - 0}{0 + 2} = \frac{1}{2}$ ,因为 $\overrightarrow{AB}//\overrightarrow{FP}$ ,所以 ${k}_{FP} = \frac{1}{2}$ ,}
{\small 所以直线 ${FP}$ 的方程为 $y = \frac{1}{2}\left( {x + \sqrt{3}}\right)$ ,}
{\small 联立椭圆: $\left\{  {\begin{matrix} \frac{{x}^{2}}{4} + {y}^{2} = 1 \\  y = \frac{1}{2}\left( {x + \sqrt{3}}\right)  \end{matrix} \Rightarrow  {x}^{2} + {\left( x + \sqrt{3}\right) }^{2} = 4 \Rightarrow  2{x}^{2} + 2\sqrt{3}x - 1 = 0}\right.$ ,}
{\small 所以 $x = \frac{\sqrt{5} - \sqrt{3}}{2}$ 或 $x = \frac{-\sqrt{3} - \sqrt{5}}{2}$ ,}
{\small 又因为点 $P$ 位于 $y$ 轴右侧,所以 $P$ 的横坐标为 $x = \frac{\sqrt{5} - \sqrt{3}}{2}$ .}
{\small \textbf{【小问 3 详解】}}
{\small 设 $C\left( {{x}_{1},{y}_{1}}\right) , D\left( {{x}_{2},{y}_{2}}\right)$ ,联立椭圆}
{\small $\left\{  {\begin{array}{l} x - y + m = 0 \\  \frac{{x}^{2}}{4} + {y}^{2} = 1 \end{array} \Rightarrow  \frac{{\left( y - m\right) }^{2}}{4} + {y}^{2} = 1 \Rightarrow  5{y}^{2} - {2my} + {m}^{2} - 4 = 0,}\right.$}
{\small 先确定有两个交点,即 $\Delta  = {\left( -2m\right) }^{2} - 4 \times  5 \times  \left( {{m}^{2} - 4}\right)  > 0$ ,}
{\small 即 $4{m}^{2} - {20}{m}^{2} + {80} > 0 \Rightarrow   - {16}{m}^{2} + {80} > 0 \Rightarrow  {m}^{2} < 5$ ,}
{\small 所以 $- \sqrt{5} < m < \sqrt{5}$ .}
{\small 因为圆上任意一点与直径两端点连线所成的角为直角,}
{\small 而点 $A$ 在以 ${CD}$ 为直径的圆外,所以 $\angle {CAD} < {90}^{ \circ  }$ ,等价于 $\overrightarrow{AC} \cdot  \overrightarrow{AD} > 0$ ,}
{\small 由 $\overrightarrow{AC} = \left( {{x}_{1} + 2,{y}_{1}}\right) ,\overrightarrow{AD} = \left( {{x}_{2} + 2,{y}_{2}}\right)$ ,}
{\small 所以 $\overrightarrow{AC} \cdot  \overrightarrow{AD} = \left( {{x}_{1} + 2}\right) \left( {{x}_{2} + 2}\right)  + {y}_{1}{y}_{2} > 0$ ,}
{\small 即 $2{y}_{1}{y}_{2} - \left( {m - 2}\right) \left( {{y}_{1} + {y}_{2}}\right)  + {\left( m - 2\right) }^{2} > 0$ .}
{\small 将 ${y}_{1} + {y}_{2} = \frac{2m}{5},{y}_{1}{y}_{2} = \frac{{m}^{2} - 4}{5}$ 代入可得,}
{\small $2 \cdot  \frac{{m}^{2} - 4}{5} - \left( {m - 2}\right)  \cdot  \frac{2m}{5} + {\left( m - 2\right) }^{2} > 0 \Rightarrow  2\left( {{m}^{2} - 4}\right)  - {2m}\left( {m - 2}\right)  + 5{\left( m - 2\right) }^{2} > 0$ ,}
{\small $2{m}^{2} - 8 - 2{m}^{2} + {4m} + 5\left( {{m}^{2} - {4m} + 4}\right)  > 0 \Rightarrow  5{m}^{2} - {16m} + {12} > 0 \Rightarrow  \left( {{5m} - 6}\right) \left( {m - 2}\right)  > 0$ ,}
{\small 解得 $m < \frac{6}{5}$ 或 $m > 2$ ,结合 $- \sqrt{5} < m < \sqrt{5}$ ,}
{\small 所以 $m \in  \left( {-\sqrt{5},\frac{6}{5}}\right)  \cup  \left( {2,\sqrt{5}}\right)$ .}
{\small ![bo_d7obids91nqc738536r0_25_241_1031_388_283_0.jpg](images/bo_d7obids91nqc738536r0_25_241_1031_388_283_0.jpg)}
\end{solutionblock}



## 第 21 题


21. 若对于定义在 $\mathbf{R}$ 上的函数 $y = f\left( x\right)$ ,设 $I$ 是 $\mathbf{R}$ 的一个子集， ${X}_{1}$ 和 ${x}_{2}$ 是 $I$ 上任意给定的两个实数,当 ${x}_{1} \neq  {x}_{2}$ 时,恒有 $f\left( {x}_{1}\right)  \neq  f\left( {x}_{2}\right)$ ,则称函数 $y = f\left( x\right)$ 在 $I$ 上具有 “性质 $P$ ”

(1)分别判断函数 $y = {\mathrm{e}}^{x}$ 和 $y = {x}^{\frac{2}{3}}$ 是否在 $\mathbf{R}$ 上具有 “性质 $P$ ” ; (无需说明理由)

(2) 设 $I = \left\lbrack  {-2, - 1}\right\rbrack   \cup  \left\lbrack  {2,3}\right\rbrack$ ,记 $f\left( x\right)  = \left\{  \begin{array}{l}  - x + k, x \leq  k, \\  {\left( x - k\right) }^{2}, x > k \end{array}\right.$ ,若函数 $y = f\left( x\right)$ 在 $I$ 上具有 “性质 $P$ ”,求实数 $k$ 的取值范围;



\begin{solutionblock}
{\small \anslabel{【考点】：} 新定义函数、单射性质、分类讨论、函数方程}
{\small \anslabel{【答案】} (1) $y = {\mathrm{e}}^{x}$ 在 $\mathrm{R}$ 上具有 “性质 $P$ ”; $y = {x}^{\frac{2}{3}}$ 在 $\mathrm{R}$ 上不具有 “性质 $P$ ”.}
{\small (2) $\left( {-\infty , - 2\rbrack \cup \lbrack  - 1,\frac{5 - \sqrt{17}}{2}}\right)  \cup  \left( {\frac{7 - \sqrt{17}}{2},2\rbrack \cup \lbrack 3, + \infty }\right)$}
{\small (3)见解析}
{\small \anslabel{【解析】}}
{\small \anslabel{【分析】} (1) 根据 “性质 $P$ ” 的定义,判断函数 $y = {\mathrm{e}}^{x}$ 和 $y = {x}^{\frac{2}{3}}$ 在 $\mathrm{R}$ 上是否满足 “性质 $P$ ”.}
{\small (2)分情况讨论函数 $f\left( x\right)$ 在不同区间上的单调性，结合 “性质 $P$ ” 的定义确定 $k$ 的取值范围.}
{\small (3)先证明集合 $M$ 非空，再通过反证法证明集合 $M$ 是无限集或单元素集.}
{\small \textbf{【小问 1 详解】}}
{\small 因为 $y = {\mathrm{e}}^{x}$ 是严格增函数,则根据 “性质 $P$ ” 的定义可知,函数 $y = {\mathrm{e}}^{x}$ 在 $\mathrm{R}$ 上具有 “性质 $P$ ”;}
{\small 因为 ${\left( -1\right) }^{\frac{2}{3}} = {1}^{\frac{2}{3}}$ ,则函数 $y = {x}^{\frac{2}{3}}$ 在 $\mathbf{R}$ 上不具有 “性质 $P$ ”.}
{\small \textbf{【小问 2 详解】}}
{\small 当 $k \leq   - 2$ 时,此时 $f\left( x\right)  = {\left( x - k\right) }^{2}$ 在 $\left\lbrack  {-2, - 1}\right\rbrack$ 和 $\left\lbrack  {2,3}\right\rbrack$ 上单调递增,函数 $y = f\left( x\right)$ 在 $I$ 上具有“性质 $P$ ”.}
{\small 当 $k \geq  3$ 时,此时 $f\left( x\right)  =  - x + k$ 在 $\left\lbrack  {-2, - 1}\right\rbrack$ 和 $\left\lbrack  {2,3}\right\rbrack$ 上单调递减,函数 $y = f\left( x\right)$ 在 $I$ 上具有 “性质 $P$ ”.}
{\small 当 $- 2 < k <  - 1$ 时,函数 $f\left( x\right)$ 在区间 $\left\lbrack  {-2, - 1}\right\rbrack$ 不单调, $\therefore y = f\left( x\right)$ 在 $I$ 不具有 “性质 $P$ ”. 当 $- 1 \leq  k \leq  2$ 时,此时 $f\left( x\right)$ 在 $\left\lbrack  {-2, - 1}\right\rbrack$ 上单调递减,在 $\left\lbrack  {2,3}\right\rbrack$ 上单调递增,}
{\small 要使函数 $y = f\left( x\right)$ 在 $I$ 上具有 “性质 $P$ ”,}
{\small 则需满足 $f\left( {-1}\right)  > f\left( 3\right)$ 或 $f\left( {-2}\right)  < f\left( 2\right)$ ,即 $1 + k > {\left( 3 - k\right) }^{2}$ 或 $2 + k < {\left( 2 - k\right) }^{2}$ ,}
{\small 整理得 ${k}^{2} - {7k} + 8 < 0$ 或 ${k}^{2} - {5k} + 2 > 0$ ,解得 $\frac{7 - \sqrt{17}}{2} < k < \frac{7 + \sqrt{17}}{2}$ 或 $k < \frac{5 - \sqrt{17}}{2}$ 或 $k > \frac{5 + \sqrt{17}}{2}$}
{\small 又 $- 1 \leq  k \leq  2$ ,得 $- 1 \leq  k < \frac{5 - \sqrt{17}}{2}$ 或 $\frac{7 - \sqrt{17}}{2} < k \leq  2$ .}
{\small 当 $2 < k < 3$ 时，函数 $f\left( x\right)$ 在区间 $\left\lbrack  {2,3}\right\rbrack$ 不单调， $\therefore f\left( x\right)$ 在 $I$ 上不具有“性质 $P$ ”.}
{\small 综上,实数 $k$ 的取值范围是 $\left( {-\infty , - 2\rbrack \cup \lbrack  - 1,\frac{5 - \sqrt{17}}{2}}\right)  \cup  \left( {\frac{7 - \sqrt{17}}{2},2\rbrack \cup \lbrack 3, + \infty }\right)$ .}
{\small \textbf{【小问 3 详解】}}
{\small 证明: 因为函数 $y = f\left( x\right)$ 在 $\mathbf{R}$ 上具有 “性质 $P$ ”,其图像是连续曲线, $\therefore f\left( x\right)$ 在 $\mathrm{R}$ 上单调.}
{\small 若 $f\left( x\right)$ 单调递减,假设 $M$ 中至少有 2 个元素 $a < b$ ,则 $f\left( a\right)  = a, f\left( b\right)  = b,\because f\left( x\right)$ 单调递减,}
{\small $\therefore a\langle b \Rightarrow  f\left( a\right) \rangle f\left( b\right)  \Rightarrow  a > b$ ,矛盾,因此 $M$ 最多一个元素.}
{\small 若 $f\left( x\right)$ 单调递增,假设 $M$ 中至少有 2 个元素 $a < b$ ,则 $f\left( a\right)  = a, f\left( b\right)  = b$ .}
{\small 对 $\forall x \in  \left( {a, b}\right)$ ,若 $f\left( x\right)  > x$ ,由于 $f\left( x\right)$ 单调递增, $f\left( {f\left( x\right) }\right)  > f\left( x\right)  \Rightarrow  x > f\left( x\right)$ ,矛盾若 $f\left( x\right)  < x$ ,由于 $f\left( x\right)$ 单调递增, $f\left( {f\left( x\right) }\right)  < f\left( x\right)  \Rightarrow  x < f\left( x\right)$ ,矛盾.}
{\small 因此对 $\forall x \in  \left( {a, b}\right)$ ,必有 $f\left( x\right)  = x$ ,即 $\left( {a, b}\right)  \subseteq  M, M$ 是无限集.}
{\small 令 $g\left( x\right)  = f\left( x\right)  - x, g\left( x\right)$ 是连续函数,若 $f\left( x\right)$ 单调递增:}
{\small 若 $f\left( x\right)  > x$ 对所有 $x$ 都成立,则 $f\left( {f\left( x\right) }\right)  > f\left( x\right)  > x$ ,与 $f\left( {f\left( x\right) }\right)  = x$ 矛盾.}
{\small 若 $f\left( x\right)  < x$ 对所有 $x$ 都成立,则 $f\left( {f\left( x\right) }\right)  < f\left( x\right)  < x$ ,与 $f\left( {f\left( x\right) }\right)  = x$ 矛盾.}
{\small 故存在 ${x}_{1} < {x}_{2}$ ,使 $\mathrm{g}\left( {x}_{1}\right)  > 0,\mathrm{\;g}\left( {x}_{2}\right)  < 0$ ,由零点存在定理, $\exists c \in  \left( {{x}_{1},{x}_{2}}\right)$ ,使 $g\left( c\right)  = 0$ ,即 $f\left( c\right)  = c, M$ 非空.}
{\small 若 $f\left( x\right)$ 单调递减:}
{\small 当 $x \rightarrow   + \infty , f\left( x\right)  \rightarrow   - \infty , x \rightarrow   - \infty , f\left( x\right)  \rightarrow   + \infty$ ,}
{\small $\therefore x \rightarrow   + \infty , g\left( x\right)  \rightarrow   - \infty , x \rightarrow   - \infty , g\left( x\right)  \rightarrow   + \infty$ ,}
{\small 由零点存在定理, $\exists c$ 使 $g\left( c\right)  = 0, M$ 非空.}
{\small 综上, $M$ 要么是单元素集,要么是无限集,命题得证.}
\end{solutionblock}



\end{document}